% !TEX program = pdflatex
% arXiv preprint draft.  Compile on Overleaf: New Project -> Upload -> this file.
% Packages are all standard; no bibliography file needed (references inlined).
\documentclass[11pt]{article}
\usepackage[margin=1in]{geometry}
\usepackage{amsmath,amssymb}
\usepackage{booktabs}
\usepackage{array}
\usepackage{hyperref}
\usepackage{xcolor}
\usepackage{times}

\newcommand{\todo}[1]{\textcolor{red}{[TODO: #1]}}

\title{One Config Value, Thirty Points:\\
Configuration Defaults and the Reproducibility of a\\
Robot Learning Benchmark}
\author{Yandong Huang\\
\small \todo{affiliation / contact}}
\date{\today}

\begin{document}
\maketitle

\begin{abstract}
Reproducing published results in robot learning is routinely difficult, yet
\emph{how much} a reproduction is off and \emph{why} usually remain anecdotal.
Taking one widely used open-source stack (LeRobot v0.6.1) and its reference task
(PushT), we use a $2\times2$ factorial design with 200 paired rollouts per cell to
decompose the gap between the published success rate (65.4\%) and what the library
yields out of the box (33.0\%). A single configuration field, \texttt{crop\_shape},
accounts for the dominant share: enabling it alone raises success from 33.0\% to
55.0\% ($+22.0$ pp, $p<0.0001$), and with the remaining official defaults the cell
reaches 63.5\%. GroupNorm together with random backbone initialization contributes
only $+2.0$ pp on its own ($p=0.36$). Three separate experiments show that training
budget is not the lever: $2\times$, $5\times$ and $6\times$ more optimization
produce no measurable gain. We further measure the evaluation noise floor with
eleven disjoint blocks of 50 initial states: the observed block-to-block spread
($\pm6.0$ pp) matches the binomial prediction ($\pm6.2$ pp; $\chi^2=9.4$, df$=10$,
$p=0.50$), so uncertainty is fully accounted for by sample size alone. Finally, the
same configuration choice is enough to move a method comparison from
indistinguishable ($p=0.51$) to decisive ($p<0.0001$), which suggests that
configurations must be reconciled before any cross-method claim is made.
Taken together, the measurements yield a cost-normalized ordering of
interventions: the two largest gains we can produce come from changes that cost no
additional compute at all, while the only compute-expensive option we test is
measurably ineffective.
\end{abstract}

\section{Introduction}

A practitioner who installs an open-source robot learning library and trains the
reference policy on the reference task expects to land near the published number.
In our case they do not: the published Diffusion Policy result on PushT is 65.4\%
\cite{lerobot-dp}, whereas the same library with its shipped defaults yields
33.0\% --- a gap of roughly thirty percentage points, with a training loss that
looks entirely healthy.

The natural first hypotheses are (i) insufficient training, (ii) insufficient data,
and (iii) a bug. We test the first directly and reject it; the second is bounded by
prior work; and the third turns out to be the wrong frame: no line of code is
incorrect. What differs is which of the library's configuration fields are set,
and the library neither reports nor warns about the consequences.

This paper makes four contributions:
\begin{enumerate}
  \item A $2\times2$ factorial decomposition of a 30-point reproduction gap at 200
        paired rollouts per cell, attributing it primarily to one configuration
        field (\S\ref{sec:factorial}).
  \item Three independent demonstrations that training budget is not the lever
        (\S\ref{sec:budget}).
  \item The first measured evaluation noise floor for this benchmark: eleven
        disjoint seed blocks, with overdispersion formally tested
        (\S\ref{sec:noise}).
  \item Evidence that the same configuration choice flips the statistical
        conclusion of a method comparison (\S\ref{sec:method}).
  \item A cost-normalized ordering of interventions and a five-step diagnostic
        protocol for failed reproductions (\S\ref{sec:where}, \S\ref{sec:protocol}).
\end{enumerate}
We do not claim that any method is better than another; we claim that the
measurement is configuration-dependent in a way that is currently unreported.

\section{Related Work}

\textbf{Benchmark construction.} RLBench \cite{rlbench} and REPLAB
\cite{replab} standardize tasks and hardware respectively. Neither studies how
the \emph{software configuration} of a policy affects measured performance.

\textbf{Reproducibility frameworks.} MIRRER \cite{mirrer} proposes a conceptual
framework for iterated, reproduced and replicated experiments and states that a
first step is ``simply reporting the resulting performance metrics from several
iterated and reproduced experiments and calculating the variance between them.''
To our knowledge this measurement has not been reported for a VLA-style policy.
RB2 \cite{rb2} quantifies a 20\% change when the same experiment is run in two
different \emph{physical} laboratories and proposes local relative ranking. Our
setting has zero physical variation --- one machine, one room, identical hardware
--- and still exhibits a larger shift.

\textbf{Statistical evaluation.} STEP \cite{step} addresses how many trials are
needed to compare two policies in the 10--50 trial regime. Such tests presuppose
that the two measurements were taken under comparable conditions. Our results
indicate this presupposition can fail even within a single library and a single
machine. Observation-contract auditing \cite{obscontract} detects, statically from
source, silent loss of state fields across observation pathways. Our mismatch is of
a different kind: no field is lost and no code path is defective, so a static
contract audit would report both configurations as valid.

\textbf{Action chunking.} The effectiveness of action chunking in imitation
learning is acknowledged to lack a complete explanation \cite{qchunk}. We provide a
cross-architecture characterization of one of its knobs, the number of actions
executed open-loop (\S\ref{sec:na}).

\section{Experimental Setup}
\label{sec:setup}

\textbf{Stack.} LeRobot v0.6.1 with four local source patches (see the
reproducibility statement). Training and evaluation run on a single consumer GPU
(RTX 4060 Laptop, 8~GiB).

\textbf{Task and data.} \texttt{lerobot/pusht}: 206 demonstrations, 25650 frames,
$96\times96$ images. Success is defined by the environment as polygon coverage
above $0.95$; the reported reward is $\mathrm{clip}(\text{coverage}/0.95, 0, 1)$,
so reward saturates at 1.0 exactly at the success threshold.

\textbf{Evaluation protocol.} 200 episodes, 300-step horizon,
\texttt{n\_action\_steps}$\,=8$ for Diffusion Policy, a fixed evaluation seed, and
22 synchronous environments. Consecutive evaluations with the same seed share the
same initial states, so all comparisons reported here are \emph{paired}; we use the
exact McNemar test on discordant pairs and Wilson intervals for marginal rates.

\textbf{Determinism.} Two library behaviours had to be neutralized before any
comparison was meaningful: cuDNN autotuning (\texttt{cudnn.benchmark}) selected
different convolution algorithms per process, producing episode-level reward
differences of $197.30$ versus $63.91$ for an identical configuration and seed;
and action sampling drew from a global RNG whose state carried over between
evaluation batches. Both are patched. With the patches, two independent evaluations
of the same checkpoint are bit-identical.

\textbf{Design.} Four training runs share dataset, batch size (64), step count
(120k), horizon, and training seed; they differ only in whether \texttt{crop\_shape}
is $[84,84]$ and whether GroupNorm is enabled together with randomly initialized
backbone weights (the two are bound in code and cannot be separated).

\section{Results}

\subsection{A single field carries the gap}
\label{sec:factorial}

\begin{table}[t]
\centering
\caption{Success rate (\%), 200 episodes per cell, paired initial states.
Wilson 95\% intervals in brackets.}
\label{tab:factorial}
\begin{tabular}{lcc}
\toprule
 & crop off & crop on \\
\midrule
GroupNorm off (ImageNet backbone) & 33.0 [26.9, 39.8] & 55.0 [48.1, 61.7] \\
GroupNorm on (random backbone)    & 28.5 [22.7, 35.1] & \textbf{63.5} [56.6, 69.9] \\
\bottomrule
\end{tabular}
\end{table}

Table~\ref{tab:factorial} reports the $2\times2$. Main effects are $+2.0$ pp for
GroupNorm with a random backbone and $+28.5$ pp for \texttt{crop\_shape}; the
interaction is $+13.0$ pp with a bootstrap 95\% interval of $[+0.0, +25.5]$,
i.e.\ marginal. Paired McNemar tests (Table~\ref{tab:paired}) confirm that both
crop contrasts are decisive while neither GroupNorm contrast is.

\begin{table}[t]
\centering
\caption{Paired exact McNemar tests over 200 shared initial states.}
\label{tab:paired}
\begin{tabular}{lccc}
\toprule
Contrast & only first succeeds & only second succeeds & $p$ \\
\midrule
crop effect, GroupNorm off & 27 & 71 & $<0.0001$ \\
crop effect, GroupNorm on  & 20 & 90 & $<0.0001$ \\
GroupNorm effect, crop off & 43 & 34 & 0.362 \\
GroupNorm effect, crop on  & 34 & 51 & 0.082 \\
default $\rightarrow$ full recipe & 20 & 81 & $<0.0001$ \\
\bottomrule
\end{tabular}
\end{table}

Two consequences follow. First, the configuration field that separates a 33.0\%
measurement from the published 65.4\% is one whose effect is invisible in the
training loss. Second, restoring \emph{any single} official value is insufficient:
crop alone is necessary and largely sufficient, but the full recipe reaches
63.5\%, statistically indistinguishable from the published figure measured over 500
episodes.

\textbf{An upstream fix that misses this field.} The library's default for the
image encoder and for GroupNorm was changed in a merged pull request in May 2026,
whose motivation was that the defaults were ``designed for the easier PushT task''
\cite{pr3202}. \texttt{crop\_shape} was not among the changed fields, yet it is the
dominant factor here.

\subsection{Training budget is not the lever}
\label{sec:budget}

\begin{table}[t]
\centering
\caption{Success rate (\%) at two training budgets, 200 episodes each.}
\label{tab:budget}
\begin{tabular}{lcc}
\toprule
Cell & 60k steps & 120k steps \\
\midrule
default            & 31.0 (62/200) & 33.0 (66/200) \\
GroupNorm only     & 26.5 (53/200) & 28.5 (57/200) \\
crop only          & 55.0 (110/200) & 55.0 (110/200) \\
\bottomrule
\end{tabular}
\end{table}

Table~\ref{tab:budget} compares two training budgets at matched evaluation. The
crop-only cell produces \emph{exactly} the same number of successes at both budgets
(110/200). Two further observations point the same way: the full-recipe cell is
non-monotonic across five checkpoints (42.0, 48.0, 62.0, 48.0, 64.0\%) and flat
within the $\pm13$ pp interval of a 50-episode evaluation; and an ACT policy
trained for $6\times$ longer (0.8M $\rightarrow$ 4.8M samples) moves from 27.0\% to
29.5\% (31 vs 36 discordant, $p=0.63$). Three model families, three budget
multipliers, no measurable gain.

\subsection{The evaluation noise floor}
\label{sec:noise}

Eleven disjoint blocks of 50 initial states, same checkpoint and identical
configuration, give success rates of 28, 32, 18, 14, 28, 28, 18, 28, 28, 30 and
30\%. The block-to-block standard deviation is $\pm6.0$ pp against a binomial
prediction of $\pm6.2$ pp (ratio $0.97$); a $\chi^2$ overdispersion test gives
$\chi^2=9.4$ with df$=10$ ($p=0.50$). There is no evidence of extra-binomial
variance: at $n=50$, uncertainty is entirely attributable to sample size. This is,
to our knowledge, the measurement called for by \cite{mirrer} but not previously
reported, and it justifies the use of binomial / Wilson intervals throughout.

\subsection{What this does to a method comparison}
\label{sec:method}

Holding the evaluation protocol fixed, we compare ACT (100k and 200k training
steps) against Diffusion Policy under two configurations. With the library
defaults, DP at 33.0\% and ACT at 29.5\% are indistinguishable (37 vs 44
discordant, $p=0.51$). With the official recipe, DP at 63.5\% dominates (13 vs 81,
$p<0.0001$). The configuration choice alone moves the comparison from a null result
to a decisive one.

We emphasize what this does \emph{not} establish. The ACT policy was not
configuration-tuned: ACT has no official PushT recipe to restore, so we cannot
exclude that a better-configured ACT would perform differently. The claim is
therefore not about the relative merit of the two methods, but about the fragility
of the comparison itself.

\subsection{Receding-horizon length across architectures}
\label{sec:na}

For ACT, the number of actions executed before replanning has a sharp interior
optimum: success is 0, 6, 4, 12, 10, \textbf{28}, 22, 4 and 0\% for
\texttt{n\_action\_steps} of 1, 2, 4, 6, 8, 10, 20, 50 and 100 respectively
(50 episodes each). For SmolVLA, tuned on the same data, the ordering is 2\% at
\texttt{n}=50, 14\% at 20, and 24\% at both 10 and 5; the 24\% vs 2\% contrast is
12:1 under a paired test ($p=0.0034$). The optimal open-loop horizon is 5--10
actions in both cases even though the prediction horizons (100 and 50) and model
sizes (80M and 450M) differ by a factor of five, which suggests the quantity is set
by the task's temporal scale rather than by the architecture.
\todo{These two sweeps use 50 episodes; re-run at 200 for a uniform protocol.}

\section{Where to Spend: A Cost-Normalized View}
\label{sec:where}

Because every quantity in this paper was measured under a fixed compute budget, we
can order the interventions by what they buy per unit of compute
(Table~\ref{tab:cost}). The ordering is stark.

\begin{table}[t]
\centering
\caption{Interventions measured in this study, ordered by cost. ``n.s.'' marks a
difference that did not reach significance. All figures are from the same task and
hardware; they are a ranking within this study, not a universal law.}
\label{tab:cost}
\begin{tabular}{p{4.6cm}cp{2.2cm}p{3.6cm}}
\toprule
Intervention & $\Delta$ success & Extra GPU-h & Nature of change \\
\midrule
\texttt{n\_action\_steps} 50$\rightarrow$10 (SmolVLA) & $+22$ pp & 0 & inference-time \\
\texttt{n\_action\_steps} 1$\rightarrow$10 (ACT)      & $+28$ pp & 0 & inference-time \\
\texttt{crop\_shape} none$\rightarrow[84,84]$         & $+22$ pp & 0 & configuration \\
episode budget 300$\rightarrow$600                     & $+32$ pp & 0 & task definition \\
\midrule
training $\times2$ (DP, 60k$\rightarrow$120k)          & $+2.0$ pp (n.s.) & $+7.3$ & optimization \\
training $\times6$ (ACT, 0.8M$\rightarrow$4.8M)        & $+2.5$ pp (n.s.) & $+10$ & optimization \\
\midrule
dataset reduced to $1/4$                               & $-24$ pp & --- & data \\
\bottomrule
\end{tabular}
\end{table}

The two largest gains that survive scrutiny are obtained by changing a value in a
configuration file or an inference-time parameter. Neither costs a single
additional GPU-hour. The one intervention that does cost compute --- longer
optimization --- produces differences that our sample sizes cannot distinguish from
zero, at two different multipliers and in two different model families.

We deliberately exclude the 300$\rightarrow$600 episode result from the ``gain''
column even though it is the largest number in the table: raising the horizon
changes the task definition, so the resulting success rate is not comparable to the
published 300-step protocol. It is nevertheless informative: a third of the failures
at the standard horizon are cases the policy could have solved given more steps,
which is itself an argument that the reported number is protocol-dependent.

\section{A Diagnostic Protocol for Failed Reproductions}
\label{sec:protocol}

The steps below are the ones that actually located our gap, ordered by the ratio of
information gained to time spent. Each addresses a specific symptom.

\begin{enumerate}
  \item \textbf{Establish determinism before trusting any comparison.} Evaluate the
        same checkpoint twice with identical settings. If the numbers differ, the
        comparison machinery is broken and no configuration study is meaningful.
        In our stack this required disabling cuDNN autotuning and re-seeding the
        action-sampling RNG once per evaluation batch.
  \item \textbf{Verify the metric, not just the number.} Confirm whether the
        reported quantity is a thresholded success indicator or a clipped continuous
        score; the two are not monotonically recoverable from one another, and a
        saturating score can rank methods differently from the success rate.
  \item \textbf{Diff configurations against the reference artifact.} Field-by-field
        comparison of the stored training configuration against the published
        checkpoint's configuration located all four discrepancies in this study.
        Visual inspection of two YAML files does not scale; a sorted diff does.
  \item \textbf{Use a factorial rather than one-factor-at-a-time.} A single-factor
        sweep would have concluded that the normalization change was useless (it
        costs 4.5 pp on its own) and stopped. The $2\times2$ revealed that the
        effect of the input crop depends on the normalization choice.
  \item \textbf{Pair the evaluations and report $n$ with an interval.} Reusing the
        same initial states across configurations multiplies the sensitivity of the
        comparison; reporting a bare success rate with $n=50$ hides a $\pm12$ pp
        interval that is larger than most claimed improvements.
\end{enumerate}

\section{Discussion}

The distinction that matters is between two kinds of mismatch. An
\emph{implementation defect} loses or reorders information: a field is dropped, a
key is resorted, a code path is skipped. Such defects are in principle detectable
without running anything, and a static contract audit \cite{obscontract} is the
right tool for them. A \emph{convention mismatch} is different in kind: both
configurations are legal, both train to a healthy loss, neither loses information,
and a static audit would pass both. What fails is the comparison against a
reference produced under the other convention.

This is why the gap we measure cannot be found by inspection. It can only be found
by knowing which reference one intends to reproduce, recovering that reference's
configuration, and measuring the difference under a paired protocol. The first two
are bookkeeping; the third is statistical. The protocol in \S\ref{sec:protocol} is
an attempt to make all three routine, and the cost table in \S\ref{sec:where} is an
argument for doing them before spending compute.

\section{Limitations}

\textbf{Single task, simulation only.} All experiments use PushT. Generalization to
other environments is untested and is the most important open question.

\textbf{One training seed.} The factorial is replicated across two training budgets
but not yet across training seeds.
\todo{a second training seed is running; report when available.}

\textbf{ACT was not configuration-tuned}, as discussed in \S\ref{sec:method}.

\textbf{The reproduction target is a single published number} for a single
checkpoint, and its own value (65.4\%) was measured with 500 episodes and a
175k-step checkpoint.

\section{Reproducibility Statement}

All training and evaluation scripts, the raw per-episode results for 65 evaluation
runs, the exported source patches, and the result tables are available at
\url{https://github.com/JacksonHYD/vla-pusht-repro-artifact}. Model weights are
archived at \url{https://huggingface.co/JacksonHYD/pusht-repro-checkpoints}. The
four source patches (deterministic convolution, per-batch RNG seeding, state-less
ACT, per-episode step recording) are provided as a \texttt{.patch} file against
LeRobot commit \texttt{7e241bd6}; without them the numbers below are not
reproducible.

\begin{thebibliography}{99}

\bibitem{lerobot-dp}
LeRobot. \emph{diffusion\_pusht} model card, huggingface.co/lerobot/diffusion\_pusht.
Reports 65.4\% success over 500 episodes at the 175k-step checkpoint.

\bibitem{lerobot}
LeRobot: An Open-Source Library for End-to-End Robot Learning. arXiv, 2026.

\bibitem{rlbench}
S.~James, Z.~Ma, D.~R.~Arrojo, A.~J.~Davison.
RLBench: The Robot Learning Benchmark \& Learning Environment. arXiv:1909.12271, 2019.

\bibitem{replab}
B.~Yang, J.~Zhang, V.~Pong, S.~Levine, D.~Jayaraman.
REPLAB: A Reproducible Low-Cost Arm Benchmark Platform for Robotic Learning.
ICRA 2019. arXiv:1905.07447.

\bibitem{mirrer}
A.~Norton, B.~Flynn.
Towards Using Multiple Iterated, Reproduced, and Replicated Experiments with Robots
(MIRRER) for Evaluation and Benchmarking. arXiv:2408.04736, 2024.

\bibitem{rb2}
S.~Dasari et al.
RB2: Robotic Manipulation Benchmarking with a Twist. arXiv:2203.08098, 2022.

\bibitem{step}
D.~Snyder et al.
Is Your Imitation Learning Policy Better than Mine? Policy Comparison with
Near-Optimal Stopping. arXiv:2503.10966, 2025.

\bibitem{obscontract}
Observation-contract auditing for valid modality comparison in robot manipulation
benchmarks. \emph{Scientific Reports}, 2026. doi:10.1038/s41598-026-76821-1.

\bibitem{qchunk}
Q.~Li, Z.~Zhou, S.~Levine.
Reinforcement Learning with Action Chunking. NeurIPS 2025. arXiv:2507.07969.

\bibitem{pr3202}
huggingface/lerobot pull request \#3202, ``feat(policy): use pretrained vision
encoder weights by default for diffusion and vqbet'', merged 2026-05-07.

\end{thebibliography}

\end{document}